% ch4_c 第4章 §4.2尾–§4.3 (书页 166–176 / p181–p191)
%JOIN%
用 Wick 定理如下求出：
\begin{align*}
\mathrm{i}\Pi^{00}(t,\pmb{k}) &= \Theta(t)\mathcal{V}^{-1}\left\langle \Big[\sum_{\pmb{q}}\big(c^{\dagger}_{\pmb{q}}c_{\pmb{q}+\pmb{k}}\big)(t),\, \sum_{\pmb{q}'}\big(c^{\dagger}_{\pmb{q}'}c_{\pmb{q}'-\pmb{k}}\big)(0)\Big] \right\rangle \tag{4.3.1}\\
&= \Theta(t)\mathcal{V}^{-1}\sum_{\pmb{q}}\left(1-n_F(\xi_{\pmb{q}+\pmb{k}})\right)n_F(\xi_{\pmb{q}})\,\mathrm{e}^{-\mathrm{i}t(\xi_{\pmb{q}+\pmb{k}}-\xi_{\pmb{q}})}\\
&\quad -\, \Theta(t)\mathcal{V}^{-1}\sum_{\pmb{q}}\left(1-n_F(\xi_{\pmb{q}})\right)n_F(\xi_{\pmb{q}+\pmb{k}})\,\mathrm{e}^{+\mathrm{i}t(\xi_{\pmb{q}}-\xi_{\pmb{q}+\pmb{k}})}
\end{align*}
%FIG{4.10}: {嵌套费米面的密度关联 $\mathrm{i}P^{00}(0,\boldsymbol{x})$ 中的“代数长程”晶体序。}

只有当 $\pmb{q}'=\pmb{q}+\pmb{k}$ 时，上述平均才不为零。在第一项中，我们在 $\pmb{q}+\pmb{k}$ 处产生一个费米子、在 $\pmb{q}$ 处湮灭一个费米子，这给出因子 $\left(1-n_F(\xi_{\pmb{q}+\pmb{k}})\right)n_F(\xi_{\pmb{q}})$。第二项具有类似的结构，其中我们在 $\pmb{q}$ 处产生一个费米子、在 $\pmb{q}+\pmb{k}$ 处湮灭一个费米子。

在 $\omega$–$\pmb{k}$ 空间中，我们有
\begin{equation*}
\Pi^{00}(\omega,\pmb{k}) = \mathcal{V}^{-1}\sum_{\pmb{q}}\left( \frac{\left(1-n_F(\xi_{\pmb{q}+\pmb{k}})\right)n_F(\xi_{\pmb{q}})}{\omega-\xi_{\pmb{q}+\pmb{k}}+\xi_{\pmb{q}}+\mathrm{i}0^{+}} - \frac{\left(1-n_F(\xi_{\pmb{q}})\right)n_F(\xi_{\pmb{q}+\pmb{k}})}{\omega+\xi_{\pmb{q}}-\xi_{\pmb{q}+\pmb{k}}+\mathrm{i}0^{+}} \right)
\end{equation*}

$\Pi^{00}(\omega,\pmb{k})$ 的虚部仅在 $(\omega,\pmb{k})$ 对应于粒子–空穴激发的能量–动量时才不为零（见图 4.11）。这一特征出现在任何二体关联函数之中，包括流关联与自旋关联。

在 $|\pmb{k}|\ll k_F$ 且 $T$ 有限的极限下，我们有
\begin{equation*}
\Pi^{00}(\omega,\pmb{k}) = -\int \frac{\mathrm{d}^d\pmb{q}}{(2\pi)^d}\, \frac{\partial n_F}{\partial\xi}\, \frac{\pmb{k}\cdot\pmb{v}}{\omega-\pmb{k}\cdot\pmb{v}+\mathrm{i}0^{+}}
\end{equation*}

其中 $\pmb{v}=\pmb{q}/m$。在 $T\to 0$ 极限下，我们有 $\partial n_F/\partial\xi=-\delta(\xi_{\pmb{k}})$，此时积分变为对费米面的积分。让我们先考虑 $\Pi^{00}(\omega,\pmb{k})$ 的虚部：
%FIG{4.11}: {(a) 阴影区域表示粒子–空穴激发的能量–动量（对 $d>1$）。$\mathrm{Im}\Pi^{00}(\omega,\boldsymbol{k})$ 仅在阴影区域内不为零。$(\omega,\boldsymbol{k})=(0,0)$ 处的边缘具有斜率 $v_F$，由 (b) 可见。}
\begin{equation*}
\mathrm{Im}\Pi^{00}(\omega,\pmb{k}) = -\frac{\omega}{v_F k}\, \frac{k_F^{d-1}}{(2\pi)^d v_F} \int \mathrm{d}\Omega\; \pi\delta\left(\frac{\omega}{v_F k}-\cos(\theta)\right)
\end{equation*}

我们发现（见图 4.12）
\begin{equation*}
\mathrm{Im}\,\Pi^{00}(\omega,\pmb{k})
= \left\{
\begin{array}{ll}
-(\omega/v_F k)\left(k_F^2/4\pi v_F\right)\Theta(v_F k-|\omega|), & d=3,\\[1ex]
-(\omega/v_F k)\left(k_F/2\pi v_F\right)\left(1/\sqrt{1-(\omega/v_F k)^2}\right)\Theta(v_F k-|\omega|), & d=2.
\end{array}
\right.
\end{equation*}

为求实部，我们注意到 $\Pi^{00}(\omega,\pmb{k})$ 是一些形如 $1/(\omega-\epsilon+\mathrm{i}0^{+})$ 的函数之和：
\begin{equation*}
\Pi^{00}(\omega,\pmb{k}) = \int \mathrm{d}\epsilon\; \frac{\mathrm{sgn}(\epsilon)A(\epsilon,\pmb{k})}{\omega-\epsilon+\mathrm{i}0^{+}}
\end{equation*}

其中 $A(\epsilon,\pmb{k})$ 是 $\Pi^{00}(\omega,\pmb{k})$ 的谱函数。谱函数与 $\Pi^{00}(\omega,\pmb{k})$ 的虚部直接相关：
\begin{equation*}
\mathrm{Im}\Pi^{00}(\omega,\pmb{k}) = -\mathrm{sgn}(\omega)\pi A(\omega,\pmb{k})
\end{equation*}
%FIG{4.12}: {固定小 $k$ 时，(a) $d=3$ 维与 (b) $d=2$ 维的 $\mathrm{Im}\Pi^{00}$。}

这样，我们可以由虚部 $\mathrm{Im}\Pi^{00}(\omega,\pmb{k})$ 算出实部。我们发现
\begin{equation*}
\Pi^{00}(\omega,\pmb{k})
= \left\{
\begin{array}{ll}
-\dfrac{k_F m}{2\pi^2}\left[1-\dfrac{\omega}{2v_F k}\ln\left|\dfrac{\omega+v_F k}{\omega-v_F k}\right|+\mathrm{i}\dfrac{\pi\omega}{2v_F k}\Theta(v_F k-|\omega|)\right], & \text{对 } d=3,\\[4ex]
-\dfrac{m}{2\pi}\left[1-\dfrac{\omega\Theta(|\omega|-v_F k)}{\sqrt{\omega^2-v_F^2k^2}}+\mathrm{i}\dfrac{\omega\Theta(v_F k-|\omega|)}{\sqrt{v_F^2k^2-\omega^2}}\right], & \text{对 } d=2,\\[4ex]
\dfrac{1}{\pi}\dfrac{v_F k^2}{(\omega+\mathrm{i}0^{+})^2-v_F^2k^2}, & \text{对 } d=1.
\end{array}
\right. \tag{4.3.2}
\end{equation*}

我们注意到，当 $\mathrm{Im}\Pi^{00}(\omega,\pmb{k})$ 在某处（比如说 $\omega_0$）出现不连续跳变时，$\mathrm{Re}\Pi^{00}(\omega,\pmb{k})$ 会在同一处出现对数发散。我们还愿指出，编时关联函数 $P^{00}(\omega,\pmb{k})$ 由同一公式给出，只是虚部多出一个因子 $\mathrm{sgn}(\omega)$。

自由费米子的压缩率决定了势的改变能诱发多大的密度改变。它由 $\chi(\pmb{k})=-\Pi^{00}(0,\pmb{k})$ 给出（在有限波矢 $\pmb{k}$ 处），且在小 $\pmb{k}$ 下与 $\pmb{k}$ 无关。由式 (4.2.12)，我们发现压缩率在所有维度都由费米能处的态密度给出：
\begin{equation*}
\chi(\pmb{k}) = N(E_F)
\end{equation*}

正如所料。由
\begin{equation*}
\sigma(\omega) = -\lim_{\pmb{k}\to 0}\mathrm{Im}\,\frac{\omega}{k^2}\,\Pi^{00}(\omega,\pmb{k})
\end{equation*}

我们也看到光学电导率（optical conductivity）为零（$\sigma(\omega)=0$），只在 $\omega=0$ 处例外。这同样在预料之中。在没有任何相互作用时，均匀的振荡电场无法激发任何粒子–空穴激发，因而不会引起耗散。
%FIG{4.13}: {$|\boldsymbol{k}|>2k_F$ 时的 $\mathrm{Im}\Pi^{00}$。(a) 中虚线与阴影区域的交截决定了 (b) 中 $\mathrm{Im}\Pi^{00}(\omega,\boldsymbol{k})$ 的值。}

当 $\omega\ll v_F k$ 时，$\Pi^{00}$ 对所有 $d>1$ 的维度都具有如下形式：
\begin{equation*}
\Pi^{00}(\omega,\pmb{k}) = -N(E_F)\left(1+\mathrm{i}C_d\frac{\omega}{v_F k}\right)
\end{equation*}

其中 $d=2$ 时 $C_2=1$，$d=3$ 时 $C_3=\pi/2$。我们看到，当 $\pmb{k}\neq 0$ 时，振荡电场会引起耗散。电导率具有形式
\begin{equation*}
\sigma(\omega,\pmb{k}) = N(E_F)C_d\,\frac{\omega^2}{v_F k^3}
\end{equation*}
（在 $\omega\ll v_F k$ 极限下）。

一般而言，当 $\omega$ 小时，$\Pi^{00}(\omega,\pmb{k})$ 是 $\pmb{k}$ 的光滑函数，但在 $\pmb{k}=0$ 与 $|\pmb{k}|=2k_F$ 附近除外（对球形费米面）。我们已经研究了 $\Pi^{00}(\omega,\pmb{k})$ 在 $\pmb{k}=0$ 附近的奇异行为；接下来我们讨论 $\Pi^{00}(\omega,\pmb{k})$ 在 $|\pmb{k}|=2k_F$ 附近的行为。让我们考虑下式的 $T=0$ 极限：
\begin{equation*}
\Pi^{00}(\omega,\pmb{k}) = \int \frac{\mathrm{d}^2\pmb{q}}{(2\pi)^d}\left( \frac{n_F(\xi_{\pmb{q}-\frac{\pmb{k}}{2}})-n_F(\xi_{\pmb{q}+\frac{\pmb{k}}{2}})}{\omega-\left(\xi_{\pmb{q}+\frac{\pmb{k}}{2}}-\xi_{\pmb{q}-\frac{\pmb{k}}{2}}\right)+\mathrm{i}0^{+}} \right)
\end{equation*}

在图 4.13、图 4.14 与图 4.15 中，浅阴影区域内 $n_F(\xi_{\pmb{q}-\frac{\pmb{k}}{2}})-n_F(\xi_{\pmb{q}+\frac{\pmb{k}}{2}})=1$，深阴影区域内 $n_F(\xi_{\pmb{q}-\frac{\pmb{k}}{2}})-n_F(\xi_{\pmb{q}+\frac{\pmb{k}}{2}})=-1$。此外，在 $y$ 轴右侧 $\xi_{\pmb{q}+\frac{\pmb{k}}{2}}-\xi_{\pmb{q}-\frac{\pmb{k}}{2}}>0$，在 $y$ 轴左侧 $\xi_{\pmb{q}+\frac{\pmb{k}}{2}}-\xi_{\pmb{q}-\frac{\pmb{k}}{2}}<0$。虚线是 $\omega=\xi_{\pmb{q}+\frac{\pmb{k}}{2}}-\xi_{\pmb{q}-\frac{\pmb{k}}{2}}$ 的位置。虚线与阴影区域的交截决定了 $\mathrm{Im}\Pi^{00}(\omega,\pmb{k})$ 的值。从图 4.13 与图 4.14 可以清楚地看出，对小 $\omega$：当 $|\pmb{k}|>2k_F$ 时 $\mathrm{Im}\Pi^{00}(\omega,\pmb{k})=0$；当 $|\pmb{k}|<2k_F$ 时 $\mathrm{Im}\Pi^{00}(\omega,\pmb{k})\propto-\omega$；而当 $|\pmb{k}|=2k_F$ 时 $\mathrm{Im}\Pi^{00}(\omega,\pmb{k})\propto-\mathrm{sgn}(\omega)|\omega|^{(d-1)/2}$。这些结果意味着，对 $d>1$，实部 $\mathrm{Re}\Pi^{00}(0,\pmb{k})$ 在 $\pmb{k}=2k_F$ 附近有限。
%FIG{4.14}: {$|\boldsymbol{k}|<2k_F$ 时的 $\mathrm{Im}\Pi^{00}$。}

然而，对嵌套费米面，当 $\pmb{k}$ 等于嵌套波矢 $\pmb{Q}$ 时，$\mathrm{Im}\Pi^{00}(\omega,\pmb{Q})$ 在 $\omega=0$ 处有一个不连续跳变（见图 4.15）。这导致压缩率在嵌套波矢 $\pmb{Q}$ 处出现对数发散：
\begin{equation*}
\chi(\pmb{Q}) \sim N(E_F)\ln\frac{E_F}{\omega}\Big|_{\omega\to 0}
\end{equation*}

在有限温度下，$\omega=0$ 处的不连续跳变被 $T$ 抹平，对数发散被 $T$ 截断：
\begin{equation*}
\chi(\pmb{Q}) \sim N(E_F)\ln\frac{E_F}{T} \tag{4.3.3}
\end{equation*}

我们已经从等时密度关联函数看到，嵌套费米面的基态在嵌套波矢 $\pmb{Q}$ 处含有一种“代数长程”晶体序。这里我们又从零频率密度关联函数看到，打开一个具有波矢 $\pmb{Q}$ 的小周期势会感生出大的密度波。这又是一个与晶体类似的性质：在晶体中，无穷小的周期势就能钉扎晶体并感生出有限的密度波。后面我们将看到，正因为嵌套费米面离晶体序如此之近，打开一个无穷小的相互作用实际上就能导致自发对称性破缺，把费米液体态变成晶体（更确切地说，称为电荷密度波（charge-density-wave，CDW）态）。

\subsection*{问题 4.3.1}

\textbf{一维自由费米子与相互作用玻色子}\quad 计算一维自由费米子系统中大 $(x,t)$ 处的 $T=0$ 编时密度关联 $\langle T\rho(x,t)\rho(0)\rangle$。（提示：利用问题 4.2.5 的结果。）你的结果应当是式 (3.6.2) 中一维相互作用玻色子密度关联的一个特例。确定式 (3.6.2) 中 $\chi$、$v$、$K_n$ 与 $C_n$ 的取值，使之重现自由费米子的编时密度关联。

\subsection*{问题 4.3.2}

对二维自由费米子系统，计算小 $(\omega,k)$ 处虚时中的零温编时密度关联 $\mathcal{P}^{00}(\omega,k)$。完成解析延拓，得到实频率的 $\Pi^{00}$。
%FIG{4.15}: {嵌套费米面在 $\boldsymbol{k}=\boldsymbol{Q}$ 处的 $\mathrm{Im}\Pi^{00}$。}

\subsection{流算符}

\begin{keypoints}
\item 具有一般色散的粒子的流算符。
\end{keypoints}

为了计算流响应函数 $\pi^{ij}$ 与 $\Pi^{ij}$（见 3.7.2 节），我们首先需要知道流算符 $\pmb{j}$ 是什么。我们从密度算符的时间导数 $\mathrm{d}\rho/\mathrm{d}t=\mathrm{d}[c^{\dagger}(\pmb{x},t)c(\pmb{x},t)]/\mathrm{d}t$ 出发，并从守恒定律 $\mathrm{d}\rho/\mathrm{d}t+\pmb{\partial}\cdot\pmb{j}=0$ 得到 $\pmb{j}$。设费米子系统由下式描述：
\begin{equation*}
H = \sum_{\langle ij\rangle}\left(t_{ij}c^{\dagger}_i c_j + h.c.\right)
\end{equation*}

我们发现 $\mathrm{i}\partial_t c=\sum_j(t_{ij}+t^{*}_{ji})c_j$，并且
\begin{equation*}
\partial_t\rho = \sum_j j_{ij}, \qquad j_{ij} = -\mathrm{i}c^{\dagger}_i(t_{ij}+t^{*}_{ji})c_j + \mathrm{i}c^{\dagger}_j(t^{*}_{ij}+t_{ji})c_i
\end{equation*}

这些结果告诉我们，在格点上，流算符由 $j_{ij}$ 给出，它描述每单位时间从格点 $j$ 流向格点 $i$ 的粒子数。这与我们要找的东西非常不同。我们要找的是一个流算符 $\pmb{j}(\pmb{x})$，它是一个矢量，且只依赖一个坐标（而非两个）。这样的流算符对格点模型而言根本不存在。然而，如果 $\pmb{A}$ 是 $\pmb{x}$ 的光滑函数，那么 $\pmb{A}$ 与流之间的耦合只能“看到”流的光滑部分。在这个极限下，即使对格点模型，我们也能找到这样一个类矢量的流算符。

诀窍在于从一个规范化的（gauged）哈密顿量出发。我们先考虑一个连续模型
\begin{equation*}
H = \int \mathrm{d}^d\pmb{x}\; c^{\dagger}(\pmb{x})\epsilon(-\mathrm{i}\partial_{\pmb{x}})c(\pmb{x})
\end{equation*}

其中 $\epsilon(\pmb{k})$ 是费米子的能谱。规范化的哈密顿量为
\begin{equation*}
H = \int \mathrm{d}^d\pmb{x}\; c^{\dagger}(\pmb{x})\epsilon(-\mathrm{i}\partial_i + A_i(\pmb{x}))c(\pmb{x}) \tag{4.3.4}
\end{equation*}

总流算符 $\pmb{J}$ 通过下式得到
\begin{equation*}
J^i(\pmb{x}) = \frac{\delta H}{\delta A_i(\pmb{x})} \tag{4.3.5}
\end{equation*}

这样得到的流算符满足流守恒定律（见问题 4.3.3）。对二次型色散 $\epsilon(\pmb{k})=\pmb{k}^2/2m$，我们有
\begin{equation*}
J^i = \frac{1}{2m}\left[c^{\dagger}(\pmb{x})(-\mathrm{i}\partial_i)c(\pmb{x}) - \left(\mathrm{i}\partial_i c^{\dagger}(\pmb{x})\right)c(\pmb{x})\right] + \frac{1}{m}A_i(\pmb{x})c^{\dagger}(\pmb{x})c(\pmb{x}) \tag{4.3.6}
\end{equation*}

第一项具有 $m^{-1}\pmb{k}\rho$ 的形式。第二项只在 $A_i\neq 0$ 时出现，具有 $m^{-1}A_i\rho$ 的形式。由于 $m^{-1}(k_i+A_i)=v_i$ 是速度，流具有 $v_i\rho$ 的形式，正如所料。

对更一般的色散，流会相当复杂，因为 $\partial_i$ 与 $A_i(\pmb{x})$ 不对易。例如，对具有 $\epsilon(k)=\gamma k^n$ 的一维系统，我们发现如下“讨厌的”形式：
\begin{equation*}
J = \gamma\sum_{m=0}^{n-1}\left[\left(\mathrm{i}\partial_x+A_x\right)^m c^{\dagger}\right]\left[\left(-\mathrm{i}\partial_x+A_x\right)^{n-1-m}c\right]
\end{equation*}

然而，在小动量极限下，我们可以忽略 $\partial_iA_j$ 项，把 $\partial_i$ 与 $A_i(\pmb{x})$ 当作对易的量来处理。在这个极限下，我们得到
\begin{align*}
J^i &= \frac{1}{2m}\left[c^{\dagger}(\pmb{x})v^i(-\mathrm{i}\partial_{\pmb{x}})c(\pmb{x}) + \left(v^i(\mathrm{i}\partial_{\pmb{x}})c^{\dagger}(\pmb{x})\right)c(\pmb{x})\right] \tag{4.3.7}\\
&\quad + \frac{1}{2}A_j(\pmb{x})\left[c^{\dagger}(\pmb{x})K^{ij}(-\mathrm{i}\partial_{\pmb{x}})c(\pmb{x}) + \left[K^{ij}(\mathrm{i}\partial_{\pmb{x}})c^{\dagger}(\pmb{x})\right]c(\pmb{x})\right] + O(\pmb{A}^2)
\end{align*}

其中
\begin{equation*}
v^i(\pmb{k}) = \frac{\partial\epsilon(\pmb{k})}{\partial k_i}, \qquad K^{ij}(\pmb{k}) = \frac{\partial^2\epsilon(\pmb{k})}{\partial k_i\partial k_j}
\end{equation*}

在动量空间中，式 (4.3.7) 可以改写为
\begin{align*}
J^i_{\pmb{q}}(t) &= \sum_{\pmb{k}} c^{\dagger}_{\pmb{k}}c_{\pmb{k}+\pmb{q}}\,\frac{v^i(\pmb{k})+v^i(\pmb{k}+\pmb{q})}{2}\\
&\quad + \mathcal{V}^{-1}\sum_{\pmb{k}\pmb{k}'} c^{\dagger}_{\pmb{k}}c_{\pmb{k}+\pmb{k}'}A^j_{\pmb{q}-\pmb{k}'}\,\frac{K^{ij}(\pmb{k})+K^{ij}(\pmb{k}+\pmb{k}')}{2} \tag{4.3.8}
\end{align*}

其中 $c(\pmb{x},t)=\sum_{\pmb{k}}\mathcal{V}^{-1/2}c_{\pmb{k}}(t)\mathrm{e}^{\mathrm{i}\pmb{k}\cdot\pmb{x}}$，$A_i(\pmb{x},t)=\mathcal{V}^{-1}\sum_{\pmb{k}}A_{i,\pmb{k}}(t)\mathrm{e}^{\mathrm{i}\pmb{k}\cdot\pmb{x}}$，$J^i(\pmb{x},t)=\mathcal{V}^{-1}\sum_{\pmb{k}}J^i_{\pmb{k}}(t)\mathrm{e}^{\mathrm{i}\pmb{k}\cdot\pmb{x}}$。当 $\pmb{A}=0$ 时，流算符变为
\begin{equation*}
j^i_{\pmb{q}}(t) = \sum_{\pmb{k}} c^{\dagger}_{\pmb{k}}(t)c_{\pmb{k}+\pmb{q}}(t)\,\frac{v^i(\pmb{k})+v^i(\pmb{k}+\pmb{q})}{2} \tag{4.3.9}
\end{equation*}

我们想强调，上式仅在小 $\pmb{q}$ 极限下成立。因此，我们也可以把流写成 $j^i_{\pmb{q}}(t)=\sum_{\pmb{k}}c^{\dagger}_{\pmb{k}}c_{\pmb{k}+\pmb{q}}v^i(\pmb{k}+(\pmb{q}/2))$，它在 $O(\pmb{q})$ 精度内与式 (4.3.9) 一致。如果我们把 $\sum_{\pmb{k}}$ 看作对布里渊区（Brillouin zone）的求和，那么式 (4.3.8) 或式 (4.3.9) 就可以看作格点模型的（近似）流算符。

\subsection*{问题 4.3.3}

\textbf{规范不变性与流守恒}
\begin{enumerate}
\item 证明规范化的哈密顿量 (4.3.4) 在规范变换 $c(\pmb{x})\to\mathrm{e}^{-\mathrm{i}\phi(\pmb{x})}c(\pmb{x})$ 与 $A_i(\pmb{x})\to A_i(\pmb{x})+\partial_i\phi(\pmb{x})$ 下不变。
\item 在一维使用规范化的哈密顿量 (4.3.4) 计算 $\partial_t(c^{\dagger}c)$，并证明由式 (4.3.5) 定义的流满足 $\partial_t\rho+\partial_x J=0$。
\end{enumerate}

\subsection*{问题 4.3.4}

求二维自由费米子系统（色散为 $\epsilon_k=\frac{k^2}{2m}+\gamma k^4$）的流算符 $J^i$。

\subsection{流关联函数}

利用式 (4.3.9)，我们现在可以计算流响应函数的第一部分 $\mathrm{i}\pi^{ij}=\langle[j^i,j^j]\rangle$（见式 (3.7.6)）。记住，密度响应函数包含来自 $(\pmb{q},\pmb{q}+\pmb{k})$ 与 $(\pmb{q}+\pmb{k},\pmb{q})$ 处粒子–空穴的两部分贡献。当我们计算流响应函数时，同样有这两部分贡献，只不过两部分都要乘上一个额外的权重因子 $\frac{1}{2}[v^i(\pmb{q})+v^i(\pmb{q}+\pmb{k})]\frac{1}{2}[v^j(\pmb{q})+v^j(\pmb{q}+\pmb{k})]$。因此
\begin{equation*}
\pi^{ij}(\omega,\pmb{k}) = \int \frac{\mathrm{d}^2\pmb{q}}{(2\pi)^d}\, \frac{\left[n_F(\xi_{\pmb{q}})-n_F(\xi_{\pmb{q}+\pmb{k}})\right]\left[v^i(\pmb{q})+v^i(\pmb{q}+\pmb{k})\right]\left[v^j(\pmb{q})+v^j(\pmb{q}+\pmb{k})\right]}{4\left[\omega-(\xi_{\pmb{q}+\pmb{k}}-\xi_{\pmb{q}})+\mathrm{i}0^{+}\right]}
\end{equation*}

在 $\pmb{k}\ll k_F$ 极限下，我们有
\begin{equation*}
\pi^{ij}(\omega,\pmb{k}) = -\int \frac{\mathrm{d}^d\pmb{q}}{(2\pi)^d}\, \frac{\partial n_F}{\partial\xi}\, \frac{\pmb{k}\cdot\pmb{v}(\pmb{q})\left[v^i(\pmb{q})+v^i(\pmb{q}+\pmb{k})\right]\left[v^j(\pmb{q})+v^j(\pmb{q}+\pmb{k})\right]}{4\left[\omega-\pmb{k}\cdot\pmb{v}(\pmb{q})+\mathrm{i}0^{+}\right]}
\end{equation*}

我们先把自己限制在 $\epsilon=\frac{k^2}{2m}$、$T=0$、$d=2$ 的情形。我们已经算出了 $\Pi^{00}$，从而也算出了 $\Pi^{\parallel}$。为计算 $\Pi^{\perp}$，我们取 $\pmb{k}=(k,0)$，于是 $\pi^{\perp}(\omega,k)=\pi^{22}(k\hat{\pmb{x}},\omega)$。我们发现（见 4.3.6.1 节）
\begin{equation*}
\pi^{\perp}(\omega,k) = -\frac{\rho}{m} + 2\frac{\rho}{m}
\left\{
\begin{array}{ll}
\dfrac{\omega^2}{v_F^2k^2} - \mathrm{i}\dfrac{\omega}{v_F k}\sqrt{1-\dfrac{\omega^2}{v_F^2k^2}}, & \text{当 } \left|\dfrac{\omega}{v_F k}\right|<1\\[4ex]
\dfrac{\omega^2}{v_F^2k^2} - \left|\dfrac{\omega}{v_F k}\right|\sqrt{\dfrac{\omega^2}{v_F^2k^2}-1}, & \text{当 } \left|\dfrac{\omega}{v_F k}\right|>1
\end{array}
\right.
\end{equation*}

对二次型色散 $\epsilon=\frac{k^2}{2m}$，当 $\pmb{A}\neq 0$ 时流具有式 (4.3.6) 的形式，这意味着 $\Pi^{\mu\nu}$ 与 $\pi^{\mu\nu}$ 通过式 (3.7.6) 相联系。我们发现
\begin{equation*}
\Pi^{\perp}(\omega,k) = \pi^{\perp}(\omega,k) + \frac{\rho}{m} = 2\frac{\rho}{m}
\left\{
\begin{array}{ll}
\dfrac{\omega^2}{v_F^2k^2} - \mathrm{i}\dfrac{\omega}{v_F k}\sqrt{1-\dfrac{\omega^2}{v_F^2k^2}}, & \text{当 } \left|\dfrac{\omega}{v_F k}\right|<1\\[4ex]
\dfrac{\omega^2}{v_F^2k^2} - \left|\dfrac{\omega}{v_F k}\right|\sqrt{\dfrac{\omega^2}{v_F^2k^2}-1}, & \text{当 } \left|\dfrac{\omega}{v_F k}\right|>1
\end{array}
\right. \tag{4.3.10}
\end{equation*}

其中用到了 $\rho=k_F^2/4\pi$。我们注意到，接触项 $\frac{\rho}{m}$ 恰好抵消了 $\pi^{\perp}$ 中的常数项（真是一大解脱！）。若没有这一抵消，自由费米子系统就会表现得像一个超导体（见式 (3.7.11)）。

在 $d=3$ 维，我们发现（见 4.3.6.2 节）
\begin{align*}
&\Pi^{\perp}(\omega,k) \tag{4.3.11}\\
&= \frac{k_F^3}{8\pi^2 m}\left(2\frac{\omega^2}{v_F^2k^2} - \frac{\omega}{v_F k}\left(1-\frac{\omega^2}{v_F^2k^2}\right)\left(\ln\left|\frac{\omega-v_F k}{\omega+v_F k}\right| + \mathrm{i}\pi\Theta(v_F k-|\omega|)\right)\right)
\end{align*}

\subsection*{问题 4.3.5}

对一维自由费米子系统精确计算 $\Pi^{\mu\nu}(\omega,k)$。画出 $\mathrm{Im}\Pi^{00}\neq 0$ 的区域，并标出 $\Pi^{00}(\omega,k)$ 变为非解析的所有边界线。指出与这些边界线相对应的粒子–空穴激发。

\subsection{电导率}

\begin{keypoints}
\item 有限的交流电导率由破坏伽利略不变性的杂质和/或相互作用引起。
\item 电导率的一种近似计算。
\end{keypoints}

如前所述，均匀电场只与质心运动耦合。因此，$\pmb{k}=0$ 的交流电导率为 $\sigma(\omega)=-\mathrm{Im}\Pi^{\perp}(\omega,0)/\omega=-\lim_{\pmb{k}\to 0}\mathrm{Im}\frac{\omega}{k^2}\Pi^{00}(\omega,\pmb{k})=0$。要有有限的交流电导率，我们需要破坏伽利略不变性。一种办法是加入杂质。另一种破坏伽利略不变性的办法是加入一个违反伽利略不变性的相互作用项。

式 (4.2.3) 中的格林函数 $\mathrm{i}G(t,\pmb{k})$ 有如下的物理含义：如果我们在 $t=0$ 时刻在 $\pmb{k}$ 态产生一个费米子，那么 $\mathrm{i}G(t,\pmb{k})$ 就是 $t$ 时刻在 $\pmb{k}$ 态找到这个费米子的振幅。对自由费米子系统，占据 $\pmb{k}$ 态的费米子哪儿也不去，因此 $|\mathrm{i}G(t,\pmb{k})|^2=1$。存在杂质时，$\pmb{k}$ 态中的费米子会因被杂质散射而消失到具有不同动量的态中去。类似地，费米子之间的相互作用会使 $\pmb{k}$ 态中的费米子衰变到一个含两个粒子与一个空穴的多粒子态；同样地，$\pmb{k}$ 态中的费米子消失到了其他态中。在某种程度上，这样的杂质/相互作用效应可以通过在费米子传播子中加入一个衰减项来模拟，如下所示：\footnote{对含杂质系统，格林函数 $G(t_1,\pmb{x}_1;t_2,\pmb{x}_2)$ 不是 $\pmb{x}_1-\pmb{x}_2$ 的函数，因为平移对称性被破坏。在这种情况下，我们甚至无法定义 $G(t,\pmb{k})$。然而，杂质平均后的格林函数只依赖 $\pmb{x}_1-\pmb{x}_2$，此时 $G(t,\pmb{k})$ 可以定义。所以，对含杂质系统，$G(t,\pmb{k})$ 应被理解为杂质平均的格林函数。}
\begin{align*}
\mathrm{i}G(t,\pmb{k}) &= +\Theta(t)\Theta(\xi_{\pmb{k}})\mathrm{e}^{-\mathrm{i}t\xi_{\pmb{k}}-t\Gamma} - \Theta(-t)\Theta(-\xi_{\pmb{k}})\mathrm{e}^{-\mathrm{i}t\xi_{\pmb{k}}-(-t)\Gamma}\Big|_{T=0}\\
\mathrm{i}G^{\beta}(t,\pmb{k}) &= +\Theta(t)\left(1-n_F(\xi_{\pmb{k}})\right)\mathrm{e}^{-\mathrm{i}t\xi_{\pmb{k}}-t\Gamma} - \Theta(-t)n_F(\xi_{\pmb{k}})\mathrm{e}^{-\mathrm{i}t\xi_{\pmb{k}}+t\Gamma}\Big|_{T>0} \tag{4.3.12}
\end{align*}

由于 $|G(t,\pmb{k})|^2$ 以 $\mathrm{e}^{-2\Gamma t}$ 衰减，衰变率为 $2\Gamma$，而 $\tau=1/2\Gamma$ 为弛豫时间。\footnote{用衰减来模拟相互作用/杂质效应，违反了费米子数守恒。}在 $\omega$–$\pmb{k}$ 空间中，式 (4.3.12) 变为
\begin{align*}
G(\omega,\pmb{k}) &= \frac{1}{\omega-\xi_{\pmb{k}}+\mathrm{i}\Gamma\,\mathrm{sgn}(\xi_{\pmb{k}})} \tag{4.3.13}\\
G^{\beta}(\omega,\pmb{k}) &= G^{\beta}_+(\omega,\pmb{k}) + G^{\beta}_-(\omega,\pmb{k}) = \frac{1-n_F(\xi_{\pmb{k}})}{\omega-\xi_{\pmb{k}}+\mathrm{i}\Gamma} + \frac{n_F(\xi_{\pmb{k}})}{\omega-\xi_{\pmb{k}}-\mathrm{i}\Gamma} \tag{4.3.14}
\end{align*}

式 (4.3.14) 使我们能够得到费米子谱函数 $\pi A_{\pm}(\omega,\pmb{k})=\pm\mathrm{Im}G^{\beta}_{\pm}(\omega,\pmb{k})$ 如下：
\begin{equation*}
\pi A_+(\omega,\pmb{k}) = \Gamma\frac{1-n_F(\xi_{\pmb{k}})}{(\omega-\xi_{\pmb{k}})^2+\Gamma^2}, \qquad \pi A_-(\omega,\pmb{k}) = \Gamma\frac{n_F(\xi_{\pmb{k}})}{(\omega-\xi_{\pmb{k}})^2+\Gamma^2}.
\end{equation*}

然而，上述 $A_{\pm}$ 不满足式 (2.2.28)。因此，我们用式 (4.3.12) 来模拟杂质/相互作用效应的做法不是自洽的。这个问题很容易修正，只需把 $A_{\pm}$ 修改为
\begin{equation*}
\pi A_+(\omega,\pmb{k}) = \Gamma\frac{1-n_F(\omega)}{(\omega-\xi_{\pmb{k}})^2+\Gamma^2}, \qquad \pi A_-(\omega,\pmb{k}) = \Gamma\frac{n_F(\omega)}{(\omega-\xi_{\pmb{k}})^2+\Gamma^2}. \tag{4.3.15}
\end{equation*}

当 $\Gamma=0^{+}$ 时，修改后的 $A_{\pm}$ 就成为自由费米子的 $A_{\pm}$。修改后的格林函数不再由式 (4.3.12) 给出，而是由式 (4.3.15) 导出的那个给出。

为计算密度响应函数 $\Theta(t)\langle[\rho(t,\pmb{q}),\allowbreak\rho(0,-\pmb{q})]\rangle$，我们需要计算（见式 (4.3.1)）$\langle[c^{\dagger}_{\pmb{k}}(t)c_{\pmb{k}+\pmb{q}},\allowbreak\, c^{\dagger}_{\pmb{k}+\pmb{q}}(0)c_{\pmb{k}}(0)]\rangle$。人们很想用 Wick 定理来计算上述关联。这里我们想指出，Wick 定理只适用于自由费米子（或玻色子）系统中的关联函数；它既不适用于相互作用的费米子，也不适用于杂质系统的平均格林函数。这里我们还是使用了 Wick 定理，但只是把它当作一种近似：
\begin{align*}
&\left\langle \left[c^{\dagger}_{\pmb{k}}(t)c_{\pmb{k}+\pmb{q}},\, c^{\dagger}_{\pmb{k}+\pmb{q}}(0)c_{\pmb{k}}(0)\right] \right\rangle\\
&\approx \left\langle c_{\pmb{k}+\pmb{q}}(t)c^{\dagger}_{\pmb{k}+\pmb{q}}(0) \right\rangle \left\langle c^{\dagger}_{\pmb{k}}(t)c_{\pmb{k}}(0) \right\rangle - \left\langle c^{\dagger}_{\pmb{k}+\pmb{q}}(0)c_{\pmb{k}+\pmb{q}}(t) \right\rangle \left\langle c_{\pmb{k}}(0)c^{\dagger}_{\pmb{k}}(t) \right\rangle\\
&= \left(\mathrm{i}G(t,\pmb{k}+\pmb{q})\right)\left(-\mathrm{i}G(-t,\pmb{k})\right) - \left(-\mathrm{i}G(-t,\pmb{k}+\pmb{q})\right)^{*}\left(\mathrm{i}G(t,\pmb{k})\right)^{*}
\end{align*}

按照从式 (4.2.15) 到式 (4.2.17) 的计算，我们发现
\begin{align*}
&\mathrm{Im}\Pi^{00}(\omega,\pmb{q}) \tag{4.3.16}\\
&= \frac{\pi}{\mathcal{V}}\int \mathrm{d}\nu\, \sum_{\pmb{k}}\left(A_-(\nu,\pmb{k}+\pmb{q})A_+(\nu-\omega,\pmb{k}) - A_+(\nu,\pmb{k}+\pmb{q})A_-(\nu-\omega,\pmb{k})\right)
\end{align*}

为计算流响应函数 $\Theta(t)\langle[j^i(t,\pmb{k}),j^j(0,-\pmb{k})]\rangle$，我们只需代入权重因子 $\frac{1}{2}[v^i(\pmb{q})+v^i(\pmb{q}+\pmb{k})]\frac{1}{2}[v^j(\pmb{q})+v^j(\pmb{q}+\pmb{k})]$。我们发现
\begin{align*}
\mathrm{Im}\Pi^{ij}(\omega,\pmb{q}) = \frac{\pi}{4\mathcal{V}}\int \mathrm{d}\nu\, \sum_{\pmb{k}} \left[v^i(\pmb{k})+v^i(\pmb{q}+\pmb{k})\right]\left[v^j(\pmb{k})+v^j(\pmb{q}+\pmb{k})\right]\times\\
\left(A_-(\nu,\pmb{k}+\pmb{q})A_+(\nu-\omega,\pmb{k}) - A_+(\nu,\pmb{k}+\pmb{q})A_-(\nu-\omega,\pmb{k})\right)
\end{align*}

令 $\pmb{q}=0$，我们得到交流电导率
\begin{align*}
\sigma^{ij}(\omega) &= -\mathrm{Im}\,\frac{\Pi^{ij}(\omega,0)}{\omega}\\
&= -\int \frac{\mathrm{d}\nu}{\pi}\, \frac{\mathrm{d}^d\pmb{k}}{(2\pi)^d}\, \frac{\Gamma^2 v^i(\pmb{k})v^j(\pmb{k})\left(n_F(\nu)-n_F(\nu-\omega)\right)}{\omega\left((\nu-\xi_{\pmb{k}})^2+\Gamma^2\right)\left((\nu-\omega-\xi_{\pmb{k}})^2+\Gamma^2\right)}
\end{align*}
